\subsection*{لنگر}

در مجموعهٔ تمرین‌های «لنگر» قصد داریم با زبان \lr{C}، یک ابزار ساده برای
اجرای کانتینر بسازیم و در هر مرحله، بخشی از آن را با استفاده از مفاهیم
درس سیستم‌عامل توسعه دهیم. در نخستین گام، زیربنای اجرای برنامه را پیاده‌سازی
می‌کنیم: دریافت یک دستور، ساخت پردازهٔ فرزند، اجرای دستور در آن و انتظار
برای پایان اجرای آن. در مراحل بعد، قابلیت‌های لازم برای جداسازی محیط اجرای
پردازه‌ها به همین پروژه افزوده خواهد شد.

\paragraph{هدف و محدودهٔ این مرحله}
هدف این تمرین، درک تفاوت ایجاد پردازه و اجرای برنامه، رابطهٔ پدر و فرزند
و نحوهٔ دریافت وضعیت پایان اجرای فرزند است. برنامهٔ این مرحله، دستور را
در محیط فعلی لینوکس اجرا می‌کند و هنوز محیط ایزولهٔ کانتینر فراهم نمی‌کند.
برای انجام این تمرین صرفا نیاز به یک کرنل لینوکس که به صورت مستقیم یا \lr{Virtual Machine} بالا اورده شده باشد دارید.

\paragraph{دریافت فایل‌های تمرین}
فایل‌های اولیه، راهنمای ساخت و آزمون‌های عمومی این مرحله در مخزن زیر

\begin{latin}
\begin{verbatim}
https://github.com/MohammadAminKoohi/OS-Practical
\end{verbatim}
\end{latin}

\paragraph{رابط خط فرمان}
نام فایل اجرایی پروژه \lr{\texttt{langar}} است. پس از ساخت برنامه با
\lr{\texttt{make}}، دستور دلخواه باید با قالب زیر قابل اجرا باشد:

\begin{latin}
\begin{verbatim}
./langar run -- COMMAND [ARGUMENT ...]
./langar run -- /bin/echo "Hello OS"
./langar run -- printf '%s\n' 'one argument with spaces'
./langar run -- /bin/sleep 2
\end{verbatim}
\end{latin}

وجود جداکنندهٔ \lr{\texttt{--}} الزامی است. نخستین آرگومان پس از آن، نام
برنامه یا مسیر فایل اجرایی و آرگومان‌های بعدی، ورودی‌های همان برنامه هستند.
نام برنامه باید هم با مسیر مشخص و هم با جست‌وجو در متغیر محیطی
\lr{\texttt{PATH}} پشتیبانی شود.

پوستهٔ فراخوان، آرگومان‌ها را پیش از اجرای لنگر از یکدیگر جدا کرده است؛
بنابراین نباید آن‌ها را به یک رشته تبدیل و دوباره تجزیه کنید. برای نمونه،
در مثال دوم، عبارت دارای فاصله تنها یک آرگومان است. آرگومان خالی نیز باید
حفظ شود. نویسه‌هایی مانند نقطه‌ویرگول که داخل یک آرگومان قرار گرفته‌اند،
دادهٔ معمولی هستند و نباید توسط لنگر به‌عنوان دستور پوسته تفسیر شوند.

\paragraph{بخش اول: ایجاد فرزند و اجرای برنامه}
در فایل‌های اولیه، بخش پردازش خط فرمان در اختیار شما قرار می‌گیرد.
وظیفهٔ شما تکمیل تابع زیر در فایل \lr{\texttt{src/run.c}} است:

\begin{latin}
\begin{verbatim}
int run_command(char *const command_argv[]);
\end{verbatim}
\end{latin}

\begin{enumerate}
  \item با استفاده از \lr{\texttt{fork()}}، یک فرزند مستقیم برای پردازهٔ
  لنگر ایجاد کنید. در این مرحله، برای اجرای هر دستور تنها همین یک فرزند
  را ایجاد کنید؛ پردازه‌هایی که خودِ برنامهٔ اجراشده می‌سازد، مشمول این
  محدودیت نیستند.

  \item در فرزند، برنامهٔ درخواستی را با \lr{\texttt{execvp()}} و با همان
  آرگومان‌های دریافتی اجرا کنید. اگر اجرا موفق باشد، ادامهٔ کد فرزند پس از
  این فراخوان اجرا نخواهد شد.

  \item پردازهٔ پدر باید تا پایان اجرای فرزند زنده بماند و با
  \lr{\texttt{waitpid()}} وضعیت پایان آن را دریافت کند. استفاده از حلقهٔ
  انتظار فعال برای بررسی پایان فرزند مجاز نیست. اگر انتظار با خطای
  \lr{\texttt{EINTR}} قطع شد، آن را دوباره انجام دهید.

  \item ورودی استاندارد، خروجی استاندارد، خروجی خطا، متغیرهای محیطی و
  پوشهٔ کاری باید از لنگر به برنامهٔ اجراشده به ارث برسند. برای مثال،
  برنامه‌ای که از ورودی استاندارد می‌خواند، باید همچنان بتواند ورودی
  دریافت کند.
\end{enumerate}

برای تحقق هدف آموزشی این تمرین، استفاده از \lr{\texttt{system()}}،
\lr{\texttt{popen()}} و \lr{\texttt{posix\_spawn()}} مجاز نیست. همچنین
نباید به‌جای اجرای مستقیم برنامه، آن را داخل یک پوسته با
\lr{\texttt{sh -c}} اجرا کنید.

\paragraph{بخش دوم: وضعیت پایان و مدیریت خطا}
رفتار برنامه را مطابق قواعد زیر کامل کنید:

\begin{itemize}
  \item اگر فرزند به‌صورت عادی پایان یافت، کد خروج لنگر باید برابر کد
  خروج فرزند باشد؛ برای مثال، اگر برنامهٔ اجراشده مقدار
  \lr{\texttt{42}} را برگرداند، لنگر نیز باید با همان مقدار پایان یابد.

  \item اگر فرزند بر اثر یک سیگنال پایان یافت، کد خروج لنگر را برابر
  \lr{\texttt{128 + signal\_number}} قرار دهید. برای استخراج وضعیت پایان،
  مستندات ماکروهای \lr{\texttt{WIFEXITED}}، \lr{\texttt{WEXITSTATUS}}،
  \lr{\texttt{WIFSIGNALED}} و \lr{\texttt{WTERMSIG}} را مطالعه کنید.
  در این مرحله، پیاده‌سازی رسیدگی‌کنندهٔ سیگنال یا ارسال سیگنال از پدر به
  فرزند خواسته نشده است.

  \item اگر برنامه یافت نشد، یعنی \lr{\texttt{execvp()}} با خطای
  \lr{\texttt{ENOENT}} یا \lr{\texttt{ENOTDIR}} شکست خورد، کد خروج
  \lr{\texttt{127}} باشد. برای سایر خطاهای این فراخوان، مانند نداشتن
  مجوز اجرا، کد خروج \lr{\texttt{126}} باشد.

  \item در صورت شکست \lr{\texttt{fork()}} یا شکست نهایی
  \lr{\texttt{waitpid()}}، پیام خطای مناسب چاپ کنید و مقدار
  \lr{\texttt{1}} را برگردانید. برای آزمودن این حالت‌ها، منابع سیستم را
  عمداً اشباع نکنید.

  \item برای خط فرمان نامعتبر، مانند حذف نام دستور یا جداکننده،
  راهنمای کوتاه استفاده چاپ شود و کد خروج \lr{\texttt{2}} باشد.
  این رفتار در بخش پردازش خط فرمانِ فایل‌های اولیه در نظر گرفته می‌شود.

  \item پس از شکست اجرای برنامه در فرزند، پیام خطا را چاپ کنید و با
  \lr{\texttt{\_exit()}} پایان دهید. فرزند نباید وارد مسیر اجرای پدر شود.
\end{itemize}

تمام پیام‌های خطا و راهنمای استفاده باید در خروجی خطای استاندارد
\lr{\texttt{stderr}} چاپ شوند. در اجرای موفق، لنگر نباید پیام اضافی،
از جمله شناسهٔ پردازه یا پیام شروع و پایان، به خروجی دستور اضافه کند.
متن دقیق پیام‌های خطای سیستم ملاک ارزیابی نیست؛ مشخص بودن علت خطا و
رعایت کد خروج اهمیت دارد.

\paragraph{آزمون و ارزیابی}
مبنای ارزیابی عملکرد پیاده‌سازی، نتیجهٔ آزمون‌های داور خودکار روی کد
تحویلی شما است. آزمون‌ها برای قابلیت‌های مستقل، امتیازهای جداگانه دارند؛
پس درست بودن یک مثال یا نمایش اجرای موفق در ویدیو، جایگزین عبور از
آزمون‌های آن قابلیت نیست. وزن آزمون‌ها در راهنمای آزمونِ منتشرشده همراه
تمرین مشخص می‌شود. آزمون‌های ارزیابی فقط رفتارهای تصریح‌شده در این
صورت سؤال را می‌سنجند، اما ورودی آن‌ها ممکن است با مثال‌های عمومی
متفاوت باشد.

آزمون‌های عمومی، مواردی مانند اجرای دستور، حفظ آرگومان‌ها، جست‌وجوی
برنامه، وراثت محیط و ورودی و خروجی، کد خروج و مدیریت خطا را بررسی می‌کنند.
همچنین با یک برنامهٔ کمکی، بررسی می‌شود که دستور در فرزندی با شناسه‌ای
متفاوت از لنگر اجرا شود و پدر تا پایان کار آن منتظر بماند. صرفاً درست
بودن متن خروجی برای قبولی کافی نیست. استفاده از توابع خواسته‌شده و
دریافت وضعیت فرزند، در کد شما نیز بررسی خواهد شد.
گزارش و ویدیو برای بررسی توضیح شما و رعایت الزامات پیاده‌سازی استفاده
می‌شوند؛ بنابراین تحویل آن‌ها، حتی در صورت عبور از همهٔ آزمون‌ها،
الزامی است. فایل‌های آزمون یا گزارش موفقیتِ ارسالی شما، مرجع نمره‌دهی
داور نیستند؛ داور نسخهٔ آزمون‌های درس را روی برنامهٔ ساخته‌شده از کد
شما اجرا می‌کند.

شناسهٔ پردازه‌ها و ترتیب زمان‌بندی را ثابت فرض نکنید. برنامه باید با
استاندارد \lr{C11} و گزینه‌های زیر بدون هشدار ساخته شود:

\begin{latin}
\begin{verbatim}
-Wall -Wextra -Werror
\end{verbatim}
\end{latin}

\paragraph{پرسش‌های گزارش}
به پرسش‌های زیر کوتاه و با ارجاع به پیاده‌سازی خود پاسخ دهید:

\begin{enumerate}
  \item مقدار بازگشتی \lr{\texttt{fork()}} در پدر، فرزند و هنگام خطا
  چه تفاوتی دارد؟ هر مسیر در کد شما چه کاری انجام می‌دهد؟
  \item تفاوت \lr{\texttt{fork()}} و خانوادهٔ \lr{\texttt{exec}} چیست؟
  پس از اجرای موفق \lr{\texttt{execvp()}} چه اتفاقی برای برنامهٔ فرزند
  و شناسهٔ پردازهٔ آن می‌افتد؟
  \item چرا پدر باید وضعیت پایان فرزند را دریافت کند؟ اگر فرزند تمام
  شود اما پدرِ زنده این وضعیت را دریافت نکند، چه اتفاقی می‌افتد؟
  \item کدام ویژگی‌های لازم برای یک کانتینر هنوز در نسخهٔ فعلی لنگر
  وجود ندارند؟ دو مورد نام ببرید؛ پیاده‌سازی آن‌ها در این مرحله لازم نیست.
\end{enumerate}

\paragraph{ویدیوی پاسخ و توضیح کد}
یک ویدیو با تصویر خوانای صفحه ضبط کنید؛ زمان پیشنهادی آن
حدود پنج تا هشت دقیقه است. ویدیو باید مربوط به همان نسخه‌ای باشد که
تحویل می‌دهید.

\begin{itemize}
  \item مسیر اجرای پدر و فرزند، محل فراخوان‌های
  \lr{\texttt{fork()}}، \lr{\texttt{execvp()}} و
  \lr{\texttt{waitpid()}}، رسیدگی به شکست اجرا و استخراج کد خروج را
  در کد خود نشان دهید و علت انتخاب‌های مهم را توضیح دهید؛ خواندن
  خط‌به‌خط کل برنامه لازم نیست.
  \item برنامه را از روی کد بسازید و یک اجرای موفق، یک دستور ناموجود
  و اجرای دستوری با کد خروج غیرصفر را نشان دهید. در هر مورد، با
  \lr{\texttt{echo \$?}} کد خروج را بلافاصله پس از پایان لنگر نمایش
  دهید. سپس نتیجهٔ اجرای آزمون‌های عمومی را نشان دهید؛ اگر آزمونی
  ناموفق است، همان نتیجه را توضیح دهید.
\end{itemize}

علاوه بر توضیح کد، به هر دو پرسش زیر در ویدیو پاسخ دهید:

\begin{enumerate}
  \item فرض کنید برنامه‌نویسی بدون ساخت فرزند، مستقیماً در خود پردازهٔ
  لنگر، \lr{\texttt{execvp()}} را فراخوانی کرده است. چرا ممکن است
  خروجی یک دستور ساده همچنان درست به نظر برسد؟ با بررسی
  \lr{\texttt{PID}} و \lr{\texttt{PPID}} چگونه این پیاده‌سازی را
  از پیاده‌سازی خواسته‌شده تشخیص می‌دهید؟
  \item اجرای یک دستور که خودش با کد \lr{\texttt{7}} تمام می‌شود،
  چه تفاوتی با شکست \lr{\texttt{execvp()}} دارد؟ مسیر اجرای کد خود
  را در هر حالت نشان دهید و توضیح دهید چرا هنگام شکست اجرا در فرزند
  از \lr{\texttt{\_exit()}} استفاده کرده‌اید.
\end{enumerate}

\paragraph{موارد لازم برای تحویل}
پاسخ این سؤال را در یک پوشه \lr{q2}
تحویل دهید؛
کنید. پوشه‌‌ی \lr{\texttt{q2}} وباید مستقیماً
در ریشهٔ آرشیو  یا فایل فشرده شما قرار داشته باشند؛ آن‌ها را داخل یک پوشهٔ اضافی با نام
دانشجو، نام مخزن یا نام تمرین قرار ندهید. ساختار بخش مربوط به این
سؤال دقیقاً به شکل زیر باشد:

\begin{latin}
\begin{verbatim}
q2/
  Makefile
  src/
    main.c
    run.c
  include/
    runner.h
  answers.md
  test-results.txt
  video.mp4
\end{verbatim}
\end{latin}

فایل‌های اولیه را همراه تغییرات خود تحویل دهید، نه فقط فایل تغییر‌یافته
یا وصلهٔ آن را. نام فایل‌ها و کوچک و بزرگ بودن حروف آن‌ها را رعایت
کنید. فایل‌های کمکیِ کد را در صورت نیاز داخل
\lr{\texttt{src/}} و \lr{\texttt{include/}} قرار دهید.

\begin{itemize}
  \item \lr{\texttt{answers.md}}: پاسخ شماره‌گذاری‌شدهٔ پرسش‌های گزارش
  به فارسی و زمان شروع بخش‌های مهم و پاسخ دو پرسش ویدیو، مانند
  \lr{\texttt{02:10}}. در این فایل، نسخهٔ محیط اجرا و محدودیت‌های
  شناخته‌شدهٔ پیاده‌سازی را نیز کوتاه ذکر کنید.
  \item \lr{\texttt{test-results.txt}}: خروجی متنی آخرین اجرای
  آزمون‌های عمومی روی همین نسخه؛ خروجی خطا را نیز ثبت کنید.
  \item \lr{\texttt{video.mp4}}: فایل خودِ ویدیو با قالب
  \lr{MP4}؛
\end{itemize}

\textbf{نکته:} در صورت زیاد شدن حجم ویدیو می‌توانید با ابزار ffmpeg حجم آن را کاهش داده و ورژن کاهش داده شده را قرار دهید.

\paragraph{ اجرای داور خودکار}
داور پس از استخراج هر آرشیو در پوشه‌ای مستقل، از ریشهٔ آن، ساخت و
اجرای برنامه را با قالب زیر انجام می‌دهد:

\begin{latin}
\begin{verbatim}
make -C q2 clean
make -C q2 CC=gcc
./q2/langar run -- COMMAND [ARGUMENT ...]
\end{verbatim}
\end{latin}

بنابراین هدف پیش‌فرض \lr{\texttt{Makefile}} باید فایل اجرایی
\lr{\texttt{q2/langar}} را بسازد و هدف \lr{\texttt{clean}}، حتی
اگر هنوز فایلی ساخته نشده است، با موفقیت پایان یابد. متغیر
\lr{\texttt{CC}} را پشتیبانی کنید و گزینه‌های ساخت خواسته‌شده را
حفظ کنید. دستورهای ساخت باید بدون دریافت ورودی، اتصال شبکه، نصب بسته
یا دسترسی \lr{root} کار کنند و در صورت شکست، وضعیت خروج غیرصفر بدهند.
همهٔ کدهای لازم برای ساخت باید در همین پوشه موجود باشند؛ مسیر مطلقِ
متعلق به رایانهٔ خود را در ساخت یا اجرا استفاده نکنید.

پوشهٔ کاری برنامه، همان پوشه‌ای است که داور از آن لنگر را فراخوانی
می‌کند و لزوماً \lr{\texttt{q2/}} نیست. برنامه نباید هنگام شروع،
پوشهٔ کاری را به محل فایل اجرایی تغییر دهد. مقادیر شناسهٔ پردازه‌ها،
ورودی آزمون و مسیر پوشهٔ استخراج را ثابت فرض نکنید. برای هر آزمون
مهلت اجرا در نظر گرفته می‌شود و برنامه نباید پس از پایان کار فرزند
بی‌دلیل منتظر بماند.

کد و فایل‌های متنی را با کدگذاری \lr{UTF-8} تحویل دهید. فایل اجرایی،
فایل‌های میانیِ ساخت، پوشهٔ \lr{\texttt{.git}} و نسخهٔ محلی آزمون‌های
درس را در آرشیو نگذارید. داور کد را دوباره می‌سازد؛ وجود فایل اجرایی
آماده، جایگزین ساخت موفق نیست. اگر ساختار مسیرها یا رابط اجرای برنامه
با این قرارداد متفاوت باشد، آزمون‌ها امکان اجرا نخواهند داشت.
